\documentclass[11pt]{article}

\usepackage[a4paper,margin=1in]{geometry}
\usepackage[T1]{fontenc}
\usepackage[utf8]{inputenc}
\usepackage{lmodern}
\usepackage{microtype}
\usepackage{amsmath,amssymb,amsthm,mathtools}
\usepackage{booktabs}
\usepackage{xcolor}
\usepackage[colorlinks=true,linkcolor=blue!50!black,
  citecolor=blue!50!black,urlcolor=blue!50!black]{hyperref}

\newtheorem{theorem}{Theorem}[section]
\newtheorem{proposition}[theorem]{Proposition}
\newtheorem{corollary}[theorem]{Corollary}
\newtheorem{lemma}[theorem]{Lemma}
\theoremstyle{definition}
\newtheorem{definition}[theorem]{Definition}
\newtheorem{remark}[theorem]{Remark}

\newcommand{\R}{\mathbb{R}}
\newcommand{\one}{\mathbf{1}}
\newcommand{\diag}{\operatorname{diag}}
\newcommand{\conv}{\operatorname{conv}}
\newcommand{\supp}{\operatorname{supp}}
\newcommand{\osc}{\operatorname{osc}}
\newcommand{\diam}{\operatorname{diam}}
\newcommand{\dH}{d_{\mathrm H}}
\newcommand{\dd}{\,\mathrm{d}}
\newcommand{\eps}{\varepsilon}
\newcommand{\Pcal}{\mathcal{P}}
\newcommand{\Wass}{\mathcal{W}}
\newcommand{\abs}[1]{\left\lvert #1\right\rvert}
\newcommand{\norm}[1]{\left\lVert #1\right\rVert}

\title{Mean Shift through Hilbert's Projective Metric\\
  \large A Sinkhorn-Style Convergence Argument}
\author{Gabriel Peyr\'e}
\date{July 2026}

\begin{document}
\maketitle

\begin{abstract}
Blurring mean shift repeatedly replaces every particle by a kernel-weighted
average of the current cloud.  For a strictly positive kernel, each update is
therefore multiplication by a positive row-stochastic matrix, although this
matrix changes with the particle positions.  This note gives a short
convergence proof based on the same projective contraction mechanism used for
Sinkhorn scaling.  Birkhoff's theorem contracts Hilbert's metric on the
positive cone; linearizing this metric at the constant ray gives Hopf's
oscillation seminorm, whose exact Markov contraction factor is Dobrushin's
coefficient.  A uniform projective bound consequently yields geometric
consensus for discrete mean shift and exponential consensus for its
continuous-time limit.  The argument also makes precise what the analogy with
Sinkhorn does, and does not, imply.
\end{abstract}

\section{Blurring mean shift as state-dependent Markov averaging}

Classical mean shift moves a query point toward a kernel-weighted barycenter of
a \emph{fixed} data set~\cite{FukunagaHostetler1975,ComaniciuMeer2002MeanShift}.
Here we study the blurring, or self-consistent, variant: every particle moves
and the averaging measure is recomputed after each update
\cite{Chen2015BlurringMeanShift}.  This distinction is essential because only
the self-consistent dynamics is a consensus system.

Let $a_i>0$, with $\sum_{i=1}^n a_i=1$, and let
$X=(x_1,\ldots,x_n)^\top\in(\R^d)^n$.  Given a positive kernel
$K:\R^d\times\R^d\to(0,+\infty)$, define
\begin{equation}\label{eq:kernel-markov}
  (K_X)_{ij}=K(x_i,x_j),
  \qquad
  A=\diag(a_1,\ldots,a_n),
  \qquad
  P_X=\diag(K_Xa)^{-1}K_XA .
\end{equation}
The entries of $P_X$ are
\begin{equation}\label{eq:mean-shift-weights}
  (P_X)_{ij}
  =
  \frac{a_jK(x_i,x_j)}{\sum_{\ell=1}^n a_\ell K(x_i,x_\ell)}.
\end{equation}
Thus $P_X>0$ and $P_X\one=\one$.  Its $i$th row is the probability law used to
average the cloud from the viewpoint of $x_i$.

For a damping parameter $0<\tau\leq1$, discrete blurring mean shift is
\begin{equation}\label{eq:discrete-mean-shift}
  X^{k+1}
  =
  (1-\tau)X^k+\tau P_{X^k}X^k.
\end{equation}
The choice $\tau=1$ is the full barycentric update.  Its continuous-time limit
is the nonautonomous consensus system
\begin{equation}\label{eq:continuous-mean-shift}
  \dot X(t)
  =
  \bigl(P_{X(t)}-I\bigr)X(t).
\end{equation}
Every new position and every instantaneous target $P_X X$ belongs to the
convex hull of the current positions.  Consequently the convex hull is nested
along either dynamics.  In particular, if
\begin{equation*}
  C_0=\conv\{x_1^0,\ldots,x_n^0\},
\end{equation*}
then all subsequent configurations remain in the compact convex set $C_0$.

\section{From Hilbert contraction to oscillation contraction}

The projective argument is most transparent for a fixed positive matrix.  We
first recall the metric used in the convergence analysis of Sinkhorn scaling.

\begin{definition}[Hilbert metric and oscillation]
For $u,v\in\R_{++}^n$, Hilbert's projective metric is
\begin{equation}\label{eq:hilbert}
  \dH(u,v)
  =
  \max_i\log\frac{u_i}{v_i}
  -
  \min_i\log\frac{u_i}{v_i}.
\end{equation}
For $z\in\R^n$, Hopf's oscillation seminorm is
\begin{equation}\label{eq:oscillation}
  \osc(z)=\max_i z_i-\min_i z_i.
\end{equation}
Both quantities ignore the constant direction: $\dH$ is invariant under
positive rescaling, while $\osc(z+c\one)=\osc(z)$.
\end{definition}

The two geometries are related exactly by logarithmic coordinates,
\begin{equation}\label{eq:hilbert-log}
  \dH(e^z,e^{z'})=\osc(z-z'),
\end{equation}
and infinitesimally by
\begin{equation}\label{eq:hilbert-linearization}
  \dH(\one+s z,\one)
  =
  |s|\,\osc(z)+o(|s|)
  \qquad (s\to0).
\end{equation}

For a positive matrix $L$, introduce its projective cross-ratio and Birkhoff
factor
\begin{equation}\label{eq:birkhoff-factor}
  \eta(L)
  =
  \max_{i,\ell,j,r}
  \frac{L_{ij}L_{\ell r}}{L_{ir}L_{\ell j}},
  \qquad
  q(L)
  =
  \frac{\sqrt{\eta(L)}-1}{\sqrt{\eta(L)}+1}
  =
  \tanh\!\left(\frac{\log\eta(L)}{4}\right).
\end{equation}

\begin{theorem}[Birkhoff contraction]\label{thm:birkhoff}
For every positive matrix $L$ and all $u,v>0$,
\begin{equation}\label{eq:birkhoff-contraction}
  \dH(Lu,Lv)
  \leq
  q(L)\dH(u,v).
\end{equation}
\end{theorem}

This is Birkhoff's cone-contraction theorem~\cite{Birkhoff1957}.  Franklin and
Lorenz used precisely this result to obtain geometric convergence of positive
matrix scaling~\cite{FranklinLorenz1989}.

For mean shift, the state coordinates are not positive vectors and the matrix
changes at every iteration.  The correct bridge is to linearize
Theorem~\ref{thm:birkhoff} at the fixed projective ray generated by $\one$.

\begin{proposition}[Hilbert contraction linearizes to consensus contraction]
\label{prop:hilbert-to-oscillation}
Let $P>0$ be row-stochastic.  Then, for every $z\in\R^n$,
\begin{equation}\label{eq:oscillation-birkhoff}
  \osc(Pz)
  \leq
  q(P)\osc(z).
\end{equation}
\end{proposition}

\begin{proof}
Apply Theorem~\ref{thm:birkhoff} to $u_s=e^{sz}$ and $v=\one$.  Since
$P\one=\one$,
\begin{equation*}
  \dH(Pe^{sz},\one)
  \leq
  q(P)\dH(e^{sz},\one)
  =
  q(P)|s|\osc(z).
\end{equation*}
Moreover,
$Pe^{sz}=\one+sPz+O(s^2)$, so
$\log(Pe^{sz})=sPz+O(s^2)$.  Divide by $|s|$ and let $s\to0$ to obtain
\eqref{eq:oscillation-birkhoff}.
\end{proof}

The exact linear contraction factor has a classical probabilistic expression.

\begin{proposition}[Dobrushin coefficient]\label{prop:dobrushin}
For a row-stochastic matrix $P$, define
\begin{equation}\label{eq:dobrushin}
  \delta(P)
  =
  \frac12\max_{i,\ell}\sum_j|P_{ij}-P_{\ell j}|
  =
  1-min_{i,\ell}\sum_j\min(P_{ij},P_{\ell j}).
\end{equation}
Then
\begin{equation}\label{eq:dobrushin-contraction}
  \osc(Pz)\leq\delta(P)\osc(z)
  \qquad\text{for every }z\in\R^n,
\end{equation}
and $\delta(P)$ is the smallest possible constant.  If $P>0$, then
\begin{equation}\label{eq:dobrushin-birkhoff}
  \delta(P)\leq q(P)<1.
\end{equation}
\end{proposition}

\begin{proof}
For two rows $p$ and $p'$, remove their common mass
$r_j=\min(p_j,p'_j)$.  The two residual positive measures have the same total
mass
\begin{equation*}
  1-\sum_jr_j=\frac12\sum_j|p_j-p'_j|.
\end{equation*}
Their integrals against $z$ can therefore differ by at most this mass times
$\osc(z)$.  Taking the largest difference between two output coordinates gives
\eqref{eq:dobrushin-contraction}; choosing $z$ as an indicator of the positive
part of a maximizing row difference proves sharpness.  Inequality
\eqref{eq:dobrushin-birkhoff} follows from
Proposition~\ref{prop:hilbert-to-oscillation}.  This relation between
Dobrushin contraction, Hopf oscillation and the infinitesimal Hilbert metric is
developed systematically in~\cite{GaubertQu2015}.
\end{proof}

\section{Projective convergence of blurring mean shift}

We now apply the preceding fixed-matrix result to the state-dependent matrices
$P_X$.  The key observation is that row normalization and the particle weights
do not alter projective cross-ratios.  Indeed,
\begin{equation}\label{eq:cross-ratio-cancellation}
  \frac{(P_X)_{ij}(P_X)_{\ell r}}
       {(P_X)_{ir}(P_X)_{\ell j}}
  =
  \frac{K(x_i,x_j)K(x_\ell,x_r)}
       {K(x_i,x_r)K(x_\ell,x_j)}.
\end{equation}
This is the projective counterpart of the fact that multiplication by positive
diagonal matrices is an isometry of Hilbert's metric.

For a configuration $X$, write
\begin{equation}\label{eq:cloud-diameter}
  D(X)=\max_{i,\ell}\norm{x_i-x_\ell}.
\end{equation}
For every unit vector $\theta\in\R^d$, let
$z_\theta=X\theta\in\R^n$.  The elementary support-function identity
\begin{equation}\label{eq:diameter-directional}
  D(X)
  =
  \sup_{\norm{\theta}=1}\osc(z_\theta)
\end{equation}
reduces contraction in $\R^d$ to scalar oscillation contraction.

\begin{theorem}[Sinkhorn-style consensus bound]\label{thm:mean-shift-consensus}
Assume that $K$ is Lipschitz and strictly positive on $C_0\times C_0$, and
set
\begin{equation}\label{eq:kernel-bounds}
  m=\min_{C_0\times C_0}K>0,
  \qquad
  M=\max_{C_0\times C_0}K<+\infty,
  \qquad
  q=\frac{M-m}{M+m}<1.
\end{equation}
Then the following statements hold.

\begin{enumerate}
\item For the damped discrete dynamics~\eqref{eq:discrete-mean-shift}, let
\begin{equation}\label{eq:damped-factor}
  q_\tau=1-\tau(1-q)<1.
\end{equation}
Then
\begin{equation}\label{eq:discrete-diameter-rate}
  D(X^k)\leq q_\tau^kD(X^0).
\end{equation}
There exists $x_\infty\in C_0$ such that every particle converges to
$x_\infty$, with
\begin{equation}\label{eq:discrete-point-rate}
  \max_i\norm{x_i^k-x_\infty}
  \leq
  \frac{\tau D(X^0)}{1-q_\tau}q_\tau^k
  =
  \frac{D(X^0)}{1-q}q_\tau^k.
\end{equation}

\item For the continuous dynamics~\eqref{eq:continuous-mean-shift},
\begin{equation}\label{eq:continuous-diameter-rate}
  D(X(t))
  \leq
  e^{-(1-q)t}D(X(0)).
\end{equation}
There exists $x_\infty\in C_0$ such that
\begin{equation}\label{eq:continuous-point-rate}
  \max_i\norm{x_i(t)-x_\infty}
  \leq
  \frac{D(X(0))}{1-q}e^{-(1-q)t}.
\end{equation}
\end{enumerate}
\end{theorem}

\begin{proof}
The finite-dimensional vector field in~\eqref{eq:continuous-mean-shift} is
Lipschitz on $C_0^n$, and convex-hull invariance keeps every configuration
inside $C_0$; hence the continuous dynamics is globally well posed.  From
\eqref{eq:cross-ratio-cancellation} and~\eqref{eq:kernel-bounds},
\begin{equation*}
  \eta(P_X)\leq\left(\frac{M}{m}\right)^2,
  \qquad
  q(P_X)\leq\frac{M-m}{M+m}=q
\end{equation*}
uniformly along the dynamics.

For the discrete update and each direction $\theta$,
Proposition~\ref{prop:hilbert-to-oscillation} gives
\begin{align*}
  \osc(z_\theta^{k+1})
  &\leq
  (1-\tau)\osc(z_\theta^k)
  +\tau\osc(P_{X^k}z_\theta^k) \\
  &\leq
  \bigl(1-\tau+\tau q\bigr)\osc(z_\theta^k)
  =q_\tau\osc(z_\theta^k).
\end{align*}
Take the supremum over $\theta$ and iterate to prove
\eqref{eq:discrete-diameter-rate}.  Since
\begin{equation*}
  \norm{x_i^{k+1}-x_i^k}\leq\tau D(X^k),
\end{equation*}
every trajectory is Cauchy.  The vanishing diameter forces all limits to
coincide, and summing the geometric tail gives
\eqref{eq:discrete-point-rate}.

For continuous time, fix $\theta$ and write
$z(t)=X(t)\theta$.  Then
\begin{equation*}
  \dot z(t)=\bigl(P_{X(t)}-I\bigr)z(t).
\end{equation*}
The upper Dini derivative of the oscillation satisfies
\begin{equation*}
  D^+\osc(z(t))
  \leq
  \osc(P_{X(t)}z(t))-\osc(z(t))
  \leq
  -(1-q)\osc(z(t)).
\end{equation*}
Gronwall's lemma and~\eqref{eq:diameter-directional} give
\eqref{eq:continuous-diameter-rate}.  Finally,
$\norm{\dot x_i(t)}\leq D(X(t))$; integrating this bound proves convergence and
\eqref{eq:continuous-point-rate}.
\end{proof}

The theorem uses the easily read bound $m\leq K\leq M$.  The sharper intrinsic
constant is
\begin{equation}\label{eq:sharp-projective-factor}
  q_*
  =
  \sup_{X\in C_0^n}q(P_X),
\end{equation}
which can replace $q$ throughout whenever $q_*<1$.  The exact oscillation
constant is smaller still: one may replace $q(P_X)$ by $\delta(P_X)$ at each
step.  The Hilbert bound is useful because its cross-ratio formula is explicit
and invariant under all row and column scalings.

\section{Gaussian mean shift and an adaptive rate}

For the Gaussian kernel
\begin{equation}\label{eq:gaussian-kernel}
  K_\eps(x,y)=\exp\!\left(-\frac{\norm{x-y}^2}{2\eps}\right),
\end{equation}
the current convex hull has diameter $D$, so $M=1$ and
$m=e^{-D^2/(2\eps)}$.  The corresponding projective factor is
\begin{equation}\label{eq:gaussian-projective-factor}
  q(D)
  =
  \frac{1-e^{-D^2/(2\eps)}}{1+e^{-D^2/(2\eps)}}
  =
  \tanh\!\left(\frac{D^2}{4\eps}\right).
\end{equation}

\begin{corollary}[Adaptive Gaussian contraction]\label{cor:gaussian}
For the full discrete update $\tau=1$,
\begin{equation}\label{eq:gaussian-discrete-adaptive}
  D(X^{k+1})
  \leq
  \tanh\!\left(\frac{D(X^k)^2}{4\eps}\right)D(X^k)
  \leq
  \frac{D(X^k)^3}{4\eps}.
\end{equation}
For the continuous-time dynamics,
\begin{equation}\label{eq:gaussian-continuous-adaptive}
  D^+D(X(t))
  \leq
  -\left[1-\tanh\!\left(\frac{D(X(t))^2}{4\eps}\right)\right]D(X(t)).
\end{equation}
\end{corollary}

Thus strict positivity gives a global geometric rate, while contraction becomes
much stronger once the cloud is narrow relative to the bandwidth.  In
particular, the projective estimate predicts a cubic local bound for the full
discrete Gaussian update.  For a narrow kernel and a widely separated initial
cloud, however, $q(D_0)$ is extremely close to one; a long-lived multimodal
transient is therefore compatible with eventual one-point consensus.

\section{Measure-valued formulation}

The same mechanism does not depend on a finite particle representation.  For
$\alpha\in\Pcal(C_0)$, define the Markov operator
\begin{equation}\label{eq:measure-markov-operator}
  (P_\alpha f)(x)
  =
  \frac{\int_{C_0}K(x,y)f(y)\dd\alpha(y)}
       {\int_{C_0}K(x,y)\dd\alpha(y)}.
\end{equation}
It fixes constants, and the integral-operator version of Birkhoff's theorem
gives
\begin{equation}\label{eq:measure-oscillation-contraction}
  \osc_{C_0}(P_\alpha f)
  \leq
  q\,\osc_{C_0}(f),
  \qquad
  q=\frac{M-m}{M+m}.
\end{equation}
The associated mean-shift field is
\begin{equation}\label{eq:measure-mean-shift-field}
  V[\alpha](x)=P_\alpha\operatorname{Id}(x)-x,
\end{equation}
and its characteristic evolution solves
\begin{equation}\label{eq:measure-mean-shift-pde}
  \partial_t\alpha_t
  +\nabla\cdot\bigl(\alpha_tV[\alpha_t]\bigr)=0.
\end{equation}
If $K$ is Lipschitz and positive on $C_0^2$, the denominator in
\eqref{eq:measure-markov-operator} is uniformly positive, so the characteristic
flow is well posed and its convex hull remains in $C_0$.  Applying
\eqref{eq:measure-oscillation-contraction} to
$f_\theta(x)=\langle\theta,x\rangle$ gives exactly the directional-width
argument in Theorem~\ref{thm:mean-shift-consensus}.  Consequently
\begin{equation}\label{eq:measure-diameter-rate}
  \diam(\supp\alpha_t)
  \leq
  e^{-(1-q)t}\diam(\supp\alpha_0),
\end{equation}
and the characteristic trajectories converge to a common point.  Hence
$\alpha_t$ converges to a Dirac mass, with the same exponential estimate in
$\Wass_\infty$.

\section{What is shared with Sinkhorn, and what is not}

For a fixed positive matrix $K$, one Sinkhorn cycle alternates
\begin{equation}\label{eq:sinkhorn-half-steps}
  u^+=\frac{a}{Kv},
  \qquad
  v^+=\frac{b}{K^\top u^+}.
\end{equation}
Multiplication by a positive diagonal vector and entrywise inversion are
isometries of $\dH$.  Each kernel multiplication contracts by $q(K)$, so a
full Sinkhorn cycle contracts projective scaling vectors by $q(K)^2$
\cite{FranklinLorenz1989}.

Mean shift uses the same ingredients in a different order.  It first forms the
positive kernel matrix $K_X$, then row-normalizes it into the Markov matrix
$P_X$.  Diagonal normalization preserves the projective cross-ratios, while
$P_X\one=\one$ makes the constant ray a fixed point.  Linearizing Hilbert
contraction at this ray yields contraction of spatial oscillations.  The common
mechanism can therefore be summarized as follows:

\begin{center}
\begin{tabular}{@{}lll@{}}
\toprule
 & Sinkhorn & Blurring mean shift \\
\midrule
positive object & fixed Gibbs kernel $K$ & state-dependent kernel $K_X$ \\
normalization & two prescribed marginals & row sums equal to one \\
contracted quantity & scaling rays in $\dH$ & coordinate widths in $\osc$ \\
rate mechanism & Birkhoff, twice per cycle & infinitesimal Birkhoff/Dobrushin \\
limit & balanced coupling & consensus Dirac mass \\
\bottomrule
\end{tabular}
\end{center}

Three qualifications prevent overinterpreting the analogy.

\begin{enumerate}
\item The estimate contracts the diameter \emph{within one evolving cloud}.
Because $P_X$ depends on $X$, it does not by itself show that the nonlinear map
$X\mapsto P_XX$ contracts the distance between two different configurations.

\item Strict positivity is decisive.  If $K$ can vanish, then the projective
diameter can be infinite and the Birkhoff factor can be one.  Bounded-confidence
models can then converge to several disconnected clusters rather than one
consensus point~\cite{HegselmannKrause2002}.

\item The result concerns blurring mean shift.  In classical nonblurring mean
shift the reference data distribution is frozen; different initial queries may
converge to different modes even when the kernel is everywhere positive.
\end{enumerate}

The useful conceptual conclusion is therefore precise: positive-kernel mean
shift and Sinkhorn are not the same fixed-point iteration, but their linear
rates are governed by the same projective loss-of-memory mechanism.  Sinkhorn
uses Hilbert contraction globally on scaling rays; mean shift uses its
infinitesimal Hopf--Dobrushin form on spatial coordinates.

\begin{thebibliography}{99}
\small
\setlength{\itemsep}{0pt}
\setlength{\parsep}{0pt}

\bibitem{Birkhoff1957}
G.~Birkhoff.
\newblock Extensions of Jentzsch's theorem.
\newblock \emph{Transactions of the American Mathematical Society},
85(1):219--227, 1957.

\bibitem{Chen2015BlurringMeanShift}
T.-L.~Chen.
\newblock On the convergence and consistency of the blurring mean-shift
process.
\newblock \emph{Annals of the Institute of Statistical Mathematics},
67(1):157--176, 2015.

\bibitem{ComaniciuMeer2002MeanShift}
D.~Comaniciu and P.~Meer.
\newblock Mean shift: A robust approach toward feature space analysis.
\newblock \emph{IEEE Transactions on Pattern Analysis and Machine
Intelligence}, 24(5):603--619, 2002.

\bibitem{FranklinLorenz1989}
J.~Franklin and J.~Lorenz.
\newblock On the scaling of multidimensional matrices.
\newblock \emph{Linear Algebra and its Applications}, 114--115:717--735,
1989.

\bibitem{FukunagaHostetler1975}
K.~Fukunaga and L.~D.~Hostetler.
\newblock The estimation of the gradient of a density function, with
applications in pattern recognition.
\newblock \emph{IEEE Transactions on Information Theory}, 21(1):32--40,
1975.

\bibitem{GaubertQu2015}
S.~Gaubert and Z.~Qu.
\newblock Dobrushin's ergodicity coefficient for Markov operators on cones.
\newblock \emph{Integral Equations and Operator Theory}, 81(1):127--150,
2015.

\bibitem{HegselmannKrause2002}
R.~Hegselmann and U.~Krause.
\newblock Opinion dynamics and bounded confidence: Models, analysis and
simulation.
\newblock \emph{Journal of Artificial Societies and Social Simulation},
5(3), 2002.

\end{thebibliography}

\end{document}
